%\documentclass[handout,xcolor=pdftex,dvipsnames,table,mathserif]{beamer}
\documentclass[xcolor=pdftex,dvipsnames,table,mathserif]{beamer}
\usepackage{subfigure}
\usepackage{amsbsy}
\usepackage{tikz}
\usetikzlibrary{arrows}
\usepackage{amsmath,graphicx,dsfont,color}
\usepackage{amsfonts}
\usepackage{amssymb}
\usepackage{array}

\bibliographystyle{apalike}

\setbeamertemplate{bibliography item}{\insertbiblabel}
\setbeamertemplate{bibliography entry title}{}
\setbeamertemplate{bibliography entry location}{}
\setbeamertemplate{bibliography entry note}{}

\DeclareMathOperator*{\argmin}{arg\,min}
\DeclareMathOperator*{\argmax}{arg\,max}
%Definitiona

\newcommand{\x}{\mathbf{x}}
\newcommand{\X}{\mathbf{X}}
\newcommand{\W}{\mathbf{W}} %Weight
\newcommand{\bais}{\mathbf{b}}%Bais
\newcommand{\act}{\texttt{g}}%Activation
\newcommand{\loss}{L}
\newcommand{\pdata}{\hat{p}_{\texttt{data}}}
\newcommand{\nsize}{n}
\newcommand{\param}{\boldsymbol{\theta}}
\newcommand{\featmap}{\boldsymbol{\phi}}
\newcommand{\EV}{\mathbb{E}}







\usepackage{physics}

\graphicspath{{../graphics/}}

\AtBeginSection[]{
  \begin{frame}{Contents}
    \tableofcontents[currentsection, hideothersubsections]
  \end{frame}
}

\AtBeginSubsection[]{
  \begin{frame}{Contents}
    \tableofcontents[currentsection, subsectionstyle=show/shaded/hide]
  \end{frame}
}

\setbeamertemplate{footline}[frame number]{}
\setbeamertemplate{navigation symbols}{}
\setbeamertemplate{section in toc}[square]
\setbeamertemplate{items}[square]

%% For image credits on image bottom right
\usepackage[absolute,overlay]{textpos}
\setbeamercolor{framesource}{fg=gray}
\setbeamerfont{framesource}{size=\tiny}
\newcommand{\source}[1]{\begin{textblock*}{4cm}(8.7cm,8.6cm)
    \begin{beamercolorbox}[ht=0.5cm,right]{framesource}
      \usebeamerfont{framesource}\usebeamercolor[fg]{framesource} Credits: {#1}
    \end{beamercolorbox}
\end{textblock*}}



\title{Attention transformers}
\author{E. Decencière}
\date{MINES ParisTech\\
  PSL Research University\\
  Center for Mathematical Morphology
}
\titlegraphic{\includegraphics[height=1.7cm]{../graphics/logoemp}}

\useinnertheme{rounded}
\usecolortheme{rose}

%%%%%%%%%%%%%%%%%%%%%%%%%%%%%%%%%%%%%%%%%%%%%%%%%%%%%%%
%%%%%%%%%%%%%%%%%%%%%%%%%%%%%%%%%%%%%%%%%%%%%%%%%%%%%%%

\begin{document}

\frame{\titlepage}

\frame{
  \frametitle{Contents}
  \tableofcontents[hidesubsections]
}

%%%%%%%%%%%%%%%%%%%%%%%%%%%%%%%%%%%%%%%%%%%%%%%%%%
%%%%%%%%%%%%%%%%%%%%%%%%%%%%%%%%%%%%%%%%%%%%%%%%%%
\section{Introduction}


%%%%%%%%%%%%%%%%%%%%%%%%%%%%%%%%%%%%
\begin{frame}{Transformers: a new revolution in deep learning?}

\begin{itemize}
\item Transformers have brought a break-through in natural language processing
  \begin{itemize}
  \item Bidirectional Encoder Representations from Transformers (BERT, by Google~\cite{brown_language_2020})
  \item Generative Pre-trained Transformer 3 (GPT-3, by OpenAI~\cite{devlin_bert_2019}): 175 billion parameters.
  \end{itemize}
  \item They contribute to the development of new natural language processing applications (translation, voice assistants, etc.)
  \item Will they do the same in image analysis?
\end{itemize}

\end{frame}


%%%%%%%%%%%%%%%%%%%%%%%%%%%%%%%%%%%%
\begin{frame}{What are transformers?}

\begin{block}{A fist quick definition}
A transformer is a neural network architecture module that explicitly allows the network to \alert{adaptively focus its attention} on certain regions of the data.
\end{block}

\end{frame}


%%%%%%%%%%%%%%%%%%%%%%%%%%%%%%%%%%%%
%%%%%%%%%%%%%%%%%%%%%%%%%%%%%%%%%%%%
\section{Visual attention}

%%%%%%%%%%%%%%%%%%%%%%%%%%%%%%%%%%%%
\subsection{Attention in human vision}


%%%%%%%%%%%%%%%%%%%%%%%%%%%%%%%%%%%%
\begin{frame}{How do we look at an  image?}

\begin{figure}[ht]
  \centering
  \includegraphics[width=0.6\textwidth]{repin_unexpected_visitor}
  \caption{Ilya Repin, An Unexpected Visitor, 1884.}
\end{figure}

\end{frame}


%%%%%%%%%%%%%%%%%%%%%%%%%%%%%%%%%%%%
\begin{frame}{How do we look at an image?}

\begin{figure}[ht]
  \centering
  \includegraphics[width=\textwidth]{Yarbus}
    \caption{Experiments on visual attention \cite{yarbus_eye_1967}}
\end{figure}

Tasks:
\begin{itemize}
\item Age of the characters?
\item How long has the visitor been away?
\item Memorize the objects in the scene.
\end{itemize}

\end{frame}

%%%%%%%%%%%%%%%%%%%%%%%%%%%%%%%%%%%%
\begin{frame}{Information used by human visual attention}

  \begin{itemize}
  \item Bottom-up:
    \begin{itemize}
    \item local features (orientation, intensity, junctions, colour, motion, etc.)
    \item local features contrast
    \item context
    \end{itemize}
\item Top-bottom: task related
\item Construction of a single \textit{saliency map}

  \end{itemize}

\end{frame}


%%%%%%%%%%%%%%%%%%%%%%%%%%%%%%%%%%%%
\begin{frame}{Exploring the image}


\begin{columns}
  \begin{column}{.5\textwidth}
    \begin{figure}[ht]
      \centering
      \includegraphics[width=0.6\textwidth]{repin_unexpected_visitor}
    \end{figure}
  \end{column}

  \begin{column}{.5\textwidth}

    \begin{itemize}
  \item Winner-takes all! We focus on the maximum of the saliency map.
  \item Inhibition of return: We explore the following maxima, at first avoiding those that have already been inspected
  \end{itemize}

  \end{column}
\end{columns}


\end{frame}


%%%%%%%%%%%%%%%%%%%%%%%%%%%%%%%%%%%%
\begin{frame}{Why?}

  \begin{itemize}
  \item Photoreceptor cells are expensive
  \item Processing power is limited
  \item Solution: concentrate the cells in a given region and use the gaze to optimize their use
  \end{itemize}

\pause

  \begin{alertblock}{}
    \begin{itemize}
    \item The same arguments apply to artificial visual systems
    \end{itemize}
  \end{alertblock}

\end{frame}


%%%%%%%%%%%%%%%%%%%%%%%%%%%%%%%%%%%%
\subsection{Attention in image analysis}

%%%%%%%%%%%%%%%%%%%%%%%%%%%%%%%%%%%%
\begin{frame}{A classical bottom-up model}

\begin{columns}
  \begin{column}{.4\textwidth}
\begin{itemize}
\item Itti et al.~\cite{itti_model_1998} proposed a model inspired by the primate visual system.
\item It only uses low-level information.
\end{itemize}


  \end{column}

  \begin{column}{.6\textwidth}
\begin{figure}[ht]
  \centering
  \includegraphics[width=\textwidth]{visual_attention_model_itti}
\end{figure}

  \end{column}
\end{columns}



\end{frame}

%%%%%%%%%%%%%%%%%%%%%%%%%%%%%%%%%%%%
\begin{frame}{Examples \cite{itti_model_1998}}

\begin{figure}[ht]
  \centering
  \includegraphics[width=\textwidth]{saliency_maps_itti}
\end{figure}


\end{frame}

%%%%%%%%%%%%%%%%%%%%%%%%%%%%%%%%%%%%
\begin{frame}{Top-down attention models}

  \begin{itemize}
  \item These are task-dependant.
  \item Note that all detection methods can be considered as task-oriented attention methods
  \end{itemize}

  \begin{block}{Example: Face detection with the Viola-Jones method~\cite{viola_rapid_2001}}
    \begin{itemize}
    \item Define weak learners based on integrals on rectangles
    \item Select learners using AdaBoost
    \item Apply them in a hierarchical way
    \end{itemize}

\begin{figure}[ht]
  \centering
  \includegraphics[width=0.7\textwidth]{viola_jones_features}\\
  \scriptsize{Image size: $24 \times 24$ pixels}
\end{figure}


  \end{block}

\end{frame}

%%%%%%%%%%%%%%%%%%%%%%%%%%%%%%%%%%%%
\begin{frame}{Illustration~\cite{viola_rapid_2001}}

\begin{figure}[ht]
  \centering
  \includegraphics[width=0.6\textwidth]{viola_jones_example}
\end{figure}

Once attention is focused, the corresponding regions can be further analysed. Here, for identification purposes, for example.

\end{frame}


%%%%%%%%%%%%%%%%%%%%%%%%%%%%%%%%%%%%
\subsection{Attention with deep learning}




%%%%%%%%%%%%%%%%%%%%%%%%%%%%%%%%%%%%
%%%%%%%%%%%%%%%%%%%%%%%%%%%%%%%%%%%%
\section{Transformers}


%%%%%%%%%%%%%%%%%%%%%%%%%%%%%%%%%%%%
\begin{frame}{Historical land-marks}

\begin{itemize}
\item Graph transformers: \cite{lecun_gradient-based_1998}
\item Transforming auto-encoders: \cite{hinton_transforming_2011}
\end{itemize}

\end{frame}


%%%%%%%%%%%%%%%%%%%%%%%%%%%%%%%%%%%%
%%%%%%%%%%%%%%%%%%%%%%%%%%%%%%%%%%%%
\section{Conclusion}


%%%%%%%%%%%%%%%%%%%%%%%%%%%%%%%%%%%%
%%%%%%%%%%%%%%%%%%%%%%%%%%%%%%%%%%%%%%%%%%%%%%%%%%
\section*{References}

%%%%%%%%%%%%%%%%%%%%%%%%%%%%%%%%%%%%%%%%%%%%%%%%%%

\frame[allowframebreaks]{

  \scriptsize

  \frametitle{References}

  %\bibliographystyle{amsalpha}
  %\bibliographystyle{apalike}

  \bibliography{../../edf.bib}

  \normalsize

}

\end{document}
